\chapter{Chapter01 — Functions As Multilinear Polynomials}
%%% generated by the blueprint pipeline from FABL.Chapter01.FunctionsAsMultilinearPolynomials
\section{Overview}

Declarations of the Lean module \texttt{FABL.Chapter01.FunctionsAsMultilinearPolynomials}.

\section{Declarations}

\begin{definition}[Field with two elements]
\lean{FABL.𝔽₂}
\leanok
\statementsource{Analysis-of-Boolean-Functions-by-Ryan-ODonnell}{
  pdf-ffbe9c7b3306-p015-b008,
  pdf-ffbe9c7b3306-p023-b001
}
Let $\mathbb F_2$ denote the field $\mathbb Z/2\mathbb Z$ with two elements.
\end{definition}

\begin{definition}[Binary cube]
\lean{FABL.F₂Cube}
\leanok
\uses{FABL.𝔽₂}
\statementsource{lmn-constant-depth-fourier}{pdf-63759bd909d5-p003-b006}
\statementsource{Analysis-of-Boolean-Functions-by-Ryan-ODonnell}{
  pdf-ffbe9c7b3306-p020-b002,
  pdf-ffbe9c7b3306-p028-b004,
  pdf-ffbe9c7b3306-p208-b004,
  pdf-ffbe9c7b3306-p209-b005
}
For a nonnegative integer $n$, the binary cube of dimension $n$ is the set of functions
from $[n]=\{0,\ldots,n-1\}$ to $\mathbb F_2$.
\end{definition}

\begin{definition}[Integer signs]
\lean{FABL.Sign}
\statementsource{lmn-constant-depth-fourier}{pdf-63759bd909d5-p005-b002}
\statementsource{Analysis-of-Boolean-Functions-by-Ryan-ODonnell}{pdf-ffbe9c7b3306-p018-b002}
\statementsource{goldreich-approximation-method}{
  pdf-988dc9f0ae91-p008-b005,
  pdf-988dc9f0ae91-p009-b001
}
\label{FABL.Sign}
\leanok
The set of signs is the multiplicative group of units of the integers, namely $\{
-1,1\}$.
\end{definition}

\begin{definition}[Sign cube]
\lean{FABL.SignCube}
\statementsource{lmn-constant-depth-fourier}{pdf-63759bd909d5-p003-b006}
\statementsource{Analysis-of-Boolean-Functions-by-Ryan-ODonnell}{
  pdf-ffbe9c7b3306-p020-b002,
  pdf-ffbe9c7b3306-p024-b002,
  pdf-ffbe9c7b3306-p208-b004,
  pdf-ffbe9c7b3306-p209-b005
}
\statementsource{goldreich-approximation-method}{
  pdf-988dc9f0ae91-p008-b005,
  pdf-988dc9f0ae91-p009-b001
}
\label{FABL.SignCube}
\leanok
\uses{FABL.Sign}
For a nonnegative integer $n$, the sign cube of dimension $n$ is the set of functions
from $[n]=\{0,\ldots,n-1\}$ to $\{-1,1\}$.
\end{definition}

\begin{definition}[Real value of a sign]
\lean{FABL.signValue}
\statementsource{lmn-constant-depth-fourier}{pdf-63759bd909d5-p003-b006}
\statementsource{Analysis-of-Boolean-Functions-by-Ryan-ODonnell}{pdf-ffbe9c7b3306-p014-b002}
\statementsource{goldreich-approximation-method}{
  pdf-988dc9f0ae91-p008-b005,
  pdf-988dc9f0ae91-p009-b001
}
\label{FABL.signValue}
\leanok
\uses{FABL.Sign}
For a sign $s\in\{-1,1\}$, its real value is the same integer regarded as a real number.
\end{definition}

\begin{theorem}[Value of the identity sign]
\lean{FABL.signValue_one}
\label{FABL.signValue_one}
\leanok
\uses{FABL.signValue}
\difficulty{1}
The sign value of the identity element $1\in\mathbb{Z}^{\times}$ is $1$.
\end{theorem}

\begin{theorem}[Value of negative one]
\lean{FABL.signValue_neg_one}
\label{FABL.signValue_neg_one}
\leanok
\uses{FABL.signValue}
\difficulty{1}
For the sign $-1$, its real value is $-1$.
\end{theorem}

\begin{theorem}[Possible values of a sign]
\lean{FABL.signValue_eq_neg_one_or_one}
\label{FABL.signValue_eq_neg_one_or_one}
\leanok
\uses{FABL.Sign, FABL.signValue}
\difficulty{4}
For every sign $s\in\{\pm1\}$, its real value is either $-1$ or $1$.
\end{theorem}

\begin{definition}[Agreement indicator polynomial]
\lean{FABL.indicatorPolynomial}
\statementsource{Analysis-of-Boolean-Functions-by-Ryan-ODonnell}{
  pdf-ffbe9c7b3306-p021-b002,
  pdf-ffbe9c7b3306-p148-b005,
  pdf-ffbe9c7b3306-p149-b001
}
\label{FABL.indicatorPolynomial}
\leanok
\uses{FABL.signValue}
For $a,x\in\{-1,1\}^n$, define the agreement indicator polynomial by $\prod_{i=1}^n
(1+a_i x_i)/2$, where each coordinate is viewed as a real sign.
\end{definition}

\begin{theorem}[Indicator polynomial on the sign cube]
\lean{FABL.indicatorPolynomial_eq_ite}
\proofsource{Analysis-of-Boolean-Functions-by-Ryan-ODonnell}{pdf-ffbe9c7b3306-p021-b002}
\label{FABL.indicatorPolynomial_eq_ite}
\leanok
\uses{FABL.indicatorPolynomial, FABL.signValue}
\proofstep
  {FABL.Sign}
  {context}
  {Analysis-of-Boolean-Functions-by-Ryan-ODonnell}
  {pdf-ffbe9c7b3306-p014-b002}
\proofstep
  {FABL.signValue}
  {context}
  {Analysis-of-Boolean-Functions-by-Ryan-ODonnell}
  {pdf-ffbe9c7b3306-p014-b002}
\proofstep
  {FABL.signValue_one}
  {context}
  {Analysis-of-Boolean-Functions-by-Ryan-ODonnell}
  {pdf-ffbe9c7b3306-p014-b002}
\proofstep
  {FABL.signValue_neg_one}
  {context}
  {Analysis-of-Boolean-Functions-by-Ryan-ODonnell}
  {pdf-ffbe9c7b3306-p014-b002}
\proofstep
  {FABL.indicatorPolynomial}
  {direct}
  {Analysis-of-Boolean-Functions-by-Ryan-ODonnell}
  {pdf-ffbe9c7b3306-p021-b002}
\proofstep
  {FABL.monomial}
  {direct}
  {Analysis-of-Boolean-Functions-by-Ryan-ODonnell}
  {pdf-ffbe9c7b3306-p022-b001}
\proofstep
  {FABL.fourierCoeff}
  {direct}
  {Analysis-of-Boolean-Functions-by-Ryan-ODonnell}
  {pdf-ffbe9c7b3306-p025-b004}
\proofstep
  {FABL.fourierSupport}
  {context}
  {Analysis-of-Boolean-Functions-by-Ryan-ODonnell}
  {pdf-ffbe9c7b3306-p022-b002,
    pdf-ffbe9c7b3306-p018-b003}
\proofstep
  {FABL.mem_fourierSupport}
  {context}
  {Analysis-of-Boolean-Functions-by-Ryan-ODonnell}
  {pdf-ffbe9c7b3306-p022-b002}
\proofstep
  {FABL.granularityFrequencyWithin}
  {context}
  {Analysis-of-Boolean-Functions-by-Ryan-ODonnell}
  {pdf-ffbe9c7b3306-p075-b005,
    pdf-ffbe9c7b3306-p076-b001}
\proofstep
  {FABL.mem_granularityFrequencyWithin}
  {context}
  {Analysis-of-Boolean-Functions-by-Ryan-ODonnell}
  {pdf-ffbe9c7b3306-p075-b005,
    pdf-ffbe9c7b3306-p076-b001}
\difficulty{4}
For any positive integer $n$ and sign vectors $a,x\in\{ -1,1\}^n$, the polynomial
$\prod_{i=1}^n(1+a_i x_i)/2$ equals $1$ when $x=a$ and $0$ otherwise.
\end{theorem}

\begin{theorem}[Indicator polynomial reconstruction]
\lean{FABL.sum_indicatorPolynomial}
\proofsource{Analysis-of-Boolean-Functions-by-Ryan-ODonnell}{pdf-ffbe9c7b3306-p021-b002}
\label{FABL.sum_indicatorPolynomial}
\leanok
\uses{FABL.indicatorPolynomial, FABL.indicatorPolynomial_eq_ite}
\proofstep
  {FABL.Sign}
  {context}
  {Analysis-of-Boolean-Functions-by-Ryan-ODonnell}
  {pdf-ffbe9c7b3306-p014-b002}
\proofstep
  {FABL.signValue}
  {context}
  {Analysis-of-Boolean-Functions-by-Ryan-ODonnell}
  {pdf-ffbe9c7b3306-p014-b002}
\proofstep
  {FABL.signValue_one}
  {context}
  {Analysis-of-Boolean-Functions-by-Ryan-ODonnell}
  {pdf-ffbe9c7b3306-p014-b002}
\proofstep
  {FABL.signValue_neg_one}
  {context}
  {Analysis-of-Boolean-Functions-by-Ryan-ODonnell}
  {pdf-ffbe9c7b3306-p014-b002}
\proofstep
  {FABL.indicatorPolynomial}
  {direct}
  {Analysis-of-Boolean-Functions-by-Ryan-ODonnell}
  {pdf-ffbe9c7b3306-p021-b002}
\proofstep
  {FABL.indicatorPolynomial_eq_ite}
  {direct}
  {Analysis-of-Boolean-Functions-by-Ryan-ODonnell}
  {pdf-ffbe9c7b3306-p021-b002}
\proofstep
  {FABL.monomial}
  {direct}
  {Analysis-of-Boolean-Functions-by-Ryan-ODonnell}
  {pdf-ffbe9c7b3306-p022-b001}
\proofstep
  {FABL.fourierCoeff}
  {direct}
  {Analysis-of-Boolean-Functions-by-Ryan-ODonnell}
  {pdf-ffbe9c7b3306-p025-b004}
\proofstep
  {FABL.fourierSupport}
  {context}
  {Analysis-of-Boolean-Functions-by-Ryan-ODonnell}
  {pdf-ffbe9c7b3306-p022-b002}
\proofstep
  {FABL.mem_fourierSupport}
  {context}
  {Analysis-of-Boolean-Functions-by-Ryan-ODonnell}
  {pdf-ffbe9c7b3306-p022-b002}
\proofstep
  {FABL.granularityFrequencyWithin}
  {context}
  {Analysis-of-Boolean-Functions-by-Ryan-ODonnell}
  {pdf-ffbe9c7b3306-p084-b002}
\proofstep
  {FABL.mem_granularityFrequencyWithin}
  {context}
  {Analysis-of-Boolean-Functions-by-Ryan-ODonnell}
  {pdf-ffbe9c7b3306-p084-b002}
\difficulty{1}
Let $f$ be a real-valued function on the sign cube of dimension $n$. For every sign
vector $x$, its value is the sum over all sign vectors $a$ of $f(a)$ times
$\\prod_{i=1}^n (1+a_i x_i)/2$.
\end{theorem}

\begin{definition}[Sign monomial]
\lean{FABL.monomial}
\statementsource{Analysis-of-Boolean-Functions-by-Ryan-ODonnell}{
  pdf-ffbe9c7b3306-p022-b001,
  pdf-ffbe9c7b3306-p149-b001,
  pdf-ffbe9c7b3306-p208-b004,
  pdf-ffbe9c7b3306-p209-b005
}
\statementsource{goldreich-approximation-method}{
  pdf-988dc9f0ae91-p009-b001,
  pdf-988dc9f0ae91-p010-b001,
  pdf-988dc9f0ae91-p012-b003,
  pdf-988dc9f0ae91-p013-b001
}
\label{FABL.monomial}
\leanok
\uses{FABL.signValue}
For a set $S\subseteq\{1,\ldots,n\}$ and $x\in\{-1,1\}^n$, define the sign monomial
$\prod_{i\in S}x_i$.
\end{definition}

\begin{definition}[Sign monomial character]
\lean{FABL.signMonomialChar}
\statementsource{Analysis-of-Boolean-Functions-by-Ryan-ODonnell}{
  pdf-ffbe9c7b3306-p022-b004,
  pdf-ffbe9c7b3306-p023-b001,
  pdf-ffbe9c7b3306-p208-b004,
  pdf-ffbe9c7b3306-p209-b005
}
\statementsource{mossel-neeman-gaussian-noise-stability}{pdf-b53c1a5aa27d-p034-b004}
\statementsource{goldreich-approximation-method}{
  pdf-988dc9f0ae91-p009-b001,
  pdf-988dc9f0ae91-p010-b001
}
\label{FABL.signMonomialChar}
\leanok
\uses{FABL.monomial, FABL.signValue}
For a set $S\subseteq\{1,\ldots,n\}$, the sign monomial character on the additive group
of $\{-1,1\}^n$ sends $x$ to $\prod_{i\in S}x_i$.
\end{definition}

\begin{theorem}[Injectivity of Sign Monomials]
\lean{FABL.signMonomialChar_injective}
\label{FABL.signMonomialChar_injective}
\leanok
\uses{FABL.monomial, FABL.signMonomialChar, FABL.signValue}
\difficulty{5}
For every $n\in\mathbb N$, the map sending a subset $S\subseteq\{1,\ldots,n\}$ to the
character $x\mapsto\prod_{i\in S}x_i$ on the sign cube $\{-1,1\}^n$ is injective.
\end{theorem}

\begin{theorem}[Orthogonality of sign monomials]
\lean{FABL.expect_monomial_mul}
\label{FABL.expect_monomial_mul}
\leanok
\uses{FABL.monomial, FABL.signMonomialChar, FABL.signMonomialChar_injective}
\difficulty{4}
Let $n$ be a nonnegative integer and let $S,T$ be subsets of the $n$ coordinates. The
uniform average of the product of the sign monomials indexed by $S$ and $T$ is $1$ when
$S=T$ and $0$ otherwise.
\end{theorem}

\begin{definition}[Coordinate sum homomorphism]
\lean{FABL.coordinateSum}
\label{FABL.coordinateSum}
\leanok
\uses{FABL.𝔽₂}
Let $n$ be a natural number and let $S\subseteq\{0,\ldots,n-1\}$. The coordinate-sum map
sends $x\in(\mathbb F_2)^n$ to $\sum_{i\in S}x_i$ in $\mathbb F_2$; it is an additive
homomorphism.
\end{definition}

\begin{definition}[Binary sign character]
\lean{FABL.binarySign}
\statementsource{Analysis-of-Boolean-Functions-by-Ryan-ODonnell}{
  pdf-ffbe9c7b3306-p014-b002,
  pdf-ffbe9c7b3306-p023-b001
}
\statementsource{goldreich-approximation-method}{
  pdf-988dc9f0ae91-p008-b005,
  pdf-988dc9f0ae91-p009-b001
}
\label{FABL.binarySign}
\leanok
\uses{FABL.𝔽₂}
The binary sign character is the additive character on $\mathbb{Z}/2\mathbb{Z}$ that
sends the residue class of $k$ to $(-1)^k$.
\end{definition}

\begin{definition}[Coordinate parity character]
\lean{FABL.χ}
\leanok
\uses{FABL.binarySign, FABL.coordinateSum}
\statementsource{lmn-constant-depth-fourier}{pdf-63759bd909d5-p004-b002}
\statementsource{Analysis-of-Boolean-Functions-by-Ryan-ODonnell}{
  pdf-ffbe9c7b3306-p023-b001,
  pdf-ffbe9c7b3306-p091-b003,
  pdf-ffbe9c7b3306-p148-b004,
  pdf-ffbe9c7b3306-p208-b004,
  pdf-ffbe9c7b3306-p209-b005
}
\statementsource{mossel-neeman-gaussian-noise-stability}{pdf-b53c1a5aa27d-p034-b004}
For a set $S\subseteq\{1,\ldots,n\}$, define the character on
$(\mathbb{Z}/2\mathbb{Z})^n$ that sends $x$ to $(-1)^{\sum_{i\in S}x_i}$, with the
exponent interpreted modulo $2$.
\end{definition}

\begin{theorem}[Multiplicativity of binary characters]
\lean{FABL.χ_add}
\leanok
\uses{FABL.χ}
\proofsource{Analysis-of-Boolean-Functions-by-Ryan-ODonnell}{pdf-ffbe9c7b3306-p023-b001}
\proofstep
  {FABL.𝔽₂}
  {direct}
  {Analysis-of-Boolean-Functions-by-Ryan-ODonnell}
  {pdf-ffbe9c7b3306-p015-b008}
\proofstep
  {FABL.Sign}
  {context}
  {Analysis-of-Boolean-Functions-by-Ryan-ODonnell}
  {pdf-ffbe9c7b3306-p020-b001}
\proofstep
  {FABL.signValue}
  {context}
  {Analysis-of-Boolean-Functions-by-Ryan-ODonnell}
  {pdf-ffbe9c7b3306-p018-b002}
\proofstep
  {FABL.signValue_one}
  {context}
  {Analysis-of-Boolean-Functions-by-Ryan-ODonnell}
  {pdf-ffbe9c7b3306-p018-b002}
\proofstep
  {FABL.signValue_neg_one}
  {context}
  {Analysis-of-Boolean-Functions-by-Ryan-ODonnell}
  {pdf-ffbe9c7b3306-p018-b002}
\proofstep
  {FABL.monomial}
  {direct}
  {Analysis-of-Boolean-Functions-by-Ryan-ODonnell}
  {pdf-ffbe9c7b3306-p018-b008,
    pdf-ffbe9c7b3306-p022-b004}
\proofstep
  {FABL.coordinateSum}
  {context}
  {Analysis-of-Boolean-Functions-by-Ryan-ODonnell}
  {pdf-ffbe9c7b3306-p031-b002}
\proofstep
  {FABL.binarySign}
  {direct}
  {Analysis-of-Boolean-Functions-by-Ryan-ODonnell}
  {pdf-ffbe9c7b3306-p014-b002,
    pdf-ffbe9c7b3306-p022-b004}
\proofstep
  {FABL.χ}
  {direct}
  {Analysis-of-Boolean-Functions-by-Ryan-ODonnell}
  {pdf-ffbe9c7b3306-p023-b001}
\proofstep
  {FABL.fourierCoeff}
  {direct}
  {Analysis-of-Boolean-Functions-by-Ryan-ODonnell}
  {pdf-ffbe9c7b3306-p025-b004}
\proofstep
  {FABL.fourierSupport}
  {context}
  {Analysis-of-Boolean-Functions-by-Ryan-ODonnell}
  {pdf-ffbe9c7b3306-p022-b002,
    pdf-ffbe9c7b3306-p018-b003}
\proofstep
  {FABL.mem_fourierSupport}
  {context}
  {Analysis-of-Boolean-Functions-by-Ryan-ODonnell}
  {pdf-ffbe9c7b3306-p022-b002}
\proofstep
  {FABL.granularityFrequencyWithin}
  {context}
  {Analysis-of-Boolean-Functions-by-Ryan-ODonnell}
  {pdf-ffbe9c7b3306-p015-b012}
\proofstep
  {FABL.mem_granularityFrequencyWithin}
  {context}
  {Analysis-of-Boolean-Functions-by-Ryan-ODonnell}
  {pdf-ffbe9c7b3306-p015-b012}
\difficulty{1}
Let $n$ be a natural number, let $S\subseteq\mathrm{Fin}(n)$, and let $x,y\in(\mathbb
Z/2\mathbb Z)^n$. The character associated with $S$ satisfies
$\chi_S(x+y)=\chi_S(x)\chi_S(y)$.
\end{theorem}

\begin{definition}[Fourier coefficient]
\lean{FABL.fourierCoeff}
\statementsource{lmn-constant-depth-fourier}{
  pdf-63759bd909d5-p004-b003,
  pdf-63759bd909d5-p004-b004
}
\statementsource{Analysis-of-Boolean-Functions-by-Ryan-ODonnell}{
  pdf-ffbe9c7b3306-p022-b001,
  pdf-ffbe9c7b3306-p025-b004,
  pdf-ffbe9c7b3306-p091-b003,
  pdf-ffbe9c7b3306-p149-b001,
  pdf-ffbe9c7b3306-p211-b002
}
\statementsource{mossel-neeman-gaussian-noise-stability}{pdf-b53c1a5aa27d-p034-b004}
\label{FABL.fourierCoeff}
\leanok
\uses{FABL.monomial}
For a function $f:\{-1,1\}^n\to\mathbb{R}$ and a set $S\subseteq\{1,\ldots,n\}$, its
Fourier coefficient at $S$ is the uniform average
$\mathbb{E}_{x\in\{-1,1\}^n}\left[f(x)\prod_{i\in S}x_i\right]$.
\end{definition}

\begin{definition}[Fourier support]
\lean{FABL.fourierSupport}
\statementsource{lmn-constant-depth-fourier}{pdf-63759bd909d5-p004-b006}
\statementsource{Analysis-of-Boolean-Functions-by-Ryan-ODonnell}{
  pdf-ffbe9c7b3306-p018-b002,
  pdf-ffbe9c7b3306-p215-b010
}
\label{FABL.fourierSupport}
\leanok
\uses{FABL.fourierCoeff}
For a function $f:\{-1,1\}^n\to\mathbb{R}$, its Fourier support is the set of all
$S\subseteq\{1,\ldots,n\}$ for which its Fourier coefficient at $S$ is nonzero.
\end{definition}

\begin{definition}[Fourier degree]
\lean{FABL.fourierDegree}
\statementsource{lmn-constant-depth-fourier}{pdf-63759bd909d5-p004-b006}
\statementsource{Analysis-of-Boolean-Functions-by-Ryan-ODonnell}{
  pdf-ffbe9c7b3306-p028-b001,
  pdf-ffbe9c7b3306-p036-b002,
  pdf-ffbe9c7b3306-p210-b002,
  pdf-ffbe9c7b3306-p215-b009,
  pdf-ffbe9c7b3306-p215-b010
}
\label{FABL.fourierDegree}
\leanok
\uses{FABL.fourierSupport}
For a function $f:\{-1,1\}^n\to\mathbb{R}$, its Fourier degree is the largest
cardinality of a set in its Fourier support, with value $0$ when that support is empty.
\end{definition}

\begin{theorem}[Membership in the Fourier support]
\lean{FABL.mem_fourierSupport}
\label{FABL.mem_fourierSupport}
\leanok
\uses{FABL.fourierCoeff, FABL.fourierSupport}
\difficulty{1}
Let $n$ be a natural number, let $f$ be a real-valued function on the sign cube
$\\{-1,1\\}^n$, and let $S$ be a finite set of coordinates in $\text{Fin}(n)$. Then $S$
belongs to the Fourier support of $f$ exactly when the Fourier coefficient of $f$ at $S$
is nonzero.
\end{theorem}

\begin{theorem}[Fourier degree criterion]
\lean{FABL.fourierDegree_le_iff}
\proofsource{Analysis-of-Boolean-Functions-by-Ryan-ODonnell}{
  pdf-ffbe9c7b3306-p215-b010,
  pdf-ffbe9c7b3306-p216-b001,
  pdf-ffbe9c7b3306-p036-b002
}
\label{FABL.fourierDegree_le_iff}
\leanok
\uses{FABL.fourierCoeff, FABL.fourierDegree, FABL.fourierSupport, FABL.mem_fourierSupport}
\proofstep
  {FABL.Sign}
  {context}
  {Analysis-of-Boolean-Functions-by-Ryan-ODonnell}
  {pdf-ffbe9c7b3306-p014-b002,
    pdf-ffbe9c7b3306-p022-b004}
\proofstep
  {FABL.signValue}
  {context}
  {Analysis-of-Boolean-Functions-by-Ryan-ODonnell}
  {pdf-ffbe9c7b3306-p014-b002,
    pdf-ffbe9c7b3306-p022-b004}
\proofstep
  {FABL.signValue_one}
  {context}
  {Analysis-of-Boolean-Functions-by-Ryan-ODonnell}
  {pdf-ffbe9c7b3306-p014-b002}
\proofstep
  {FABL.signValue_neg_one}
  {context}
  {Analysis-of-Boolean-Functions-by-Ryan-ODonnell}
  {pdf-ffbe9c7b3306-p014-b002}
\proofstep
  {FABL.monomial}
  {direct}
  {Analysis-of-Boolean-Functions-by-Ryan-ODonnell}
  {pdf-ffbe9c7b3306-p022-b004}
\proofstep
  {FABL.fourierCoeff}
  {direct}
  {Analysis-of-Boolean-Functions-by-Ryan-ODonnell}
  {pdf-ffbe9c7b3306-p025-b004}
\proofstep
  {FABL.fourierSupport}
  {context}
  {Analysis-of-Boolean-Functions-by-Ryan-ODonnell}
  {pdf-ffbe9c7b3306-p022-b002,
    pdf-ffbe9c7b3306-p027-b008}
\proofstep
  {FABL.fourierDegree}
  {context}
  {Analysis-of-Boolean-Functions-by-Ryan-ODonnell}
  {pdf-ffbe9c7b3306-p036-b002,
    pdf-ffbe9c7b3306-p028-b002}
\proofstep
  {FABL.mem_fourierSupport}
  {context}
  {Analysis-of-Boolean-Functions-by-Ryan-ODonnell}
  {pdf-ffbe9c7b3306-p022-b002,
    pdf-ffbe9c7b3306-p027-b008}
\proofstep
  {FABL.granularityFrequencyWithin}
  {context}
  {Analysis-of-Boolean-Functions-by-Ryan-ODonnell}
  {pdf-ffbe9c7b3306-p076-b001,
    pdf-ffbe9c7b3306-p076-b003}
\proofstep
  {FABL.mem_granularityFrequencyWithin}
  {context}
  {Analysis-of-Boolean-Functions-by-Ryan-ODonnell}
  {pdf-ffbe9c7b3306-p076-b001,
    pdf-ffbe9c7b3306-p076-b003}
\difficulty{4}
For a function $f$ on the Boolean cube $\\{-1,1\\}^n$ and an integer $k\\ge 0$, its
Fourier degree is at most $k$ if and only if every Fourier coefficient indexed by a
subset $S\\subseteq[n]$ with $|S|>k$ is zero.
\end{theorem}

\begin{theorem}[Fourier degree bounded by dimension]
\lean{FABL.fourierDegree_le_dimension}
\label{FABL.fourierDegree_le_dimension}
\leanok
\uses{FABL.fourierDegree, FABL.fourierDegree_le_iff}
\difficulty{3}
For every natural number $n$ and every function $f\colon\{-1,1\}^n\to\mathbb{R}$, the
Fourier degree of $f$ is at most $n$.
\end{theorem}

\begin{definition}[Fourier granularity]
\lean{FABL.IsFourierGranular}
\statementsource{Analysis-of-Boolean-Functions-by-Ryan-ODonnell}{
  pdf-ffbe9c7b3306-p072-b003,
  pdf-ffbe9c7b3306-p089-b003
}
\label{FABL.IsFourierGranular}
\leanok
\uses{FABL.fourierCoeff}
A function $f$ on the sign cube is Fourier granular with parameter $\varepsilon$ if
every Fourier coefficient of $f$ is an integer multiple of $\varepsilon$; that is, for
each subset $S$ of coordinates, there is an integer $z$ such that $\widehat
f(S)=z\varepsilon$.
\end{definition}

\begin{definition}[Granularity indexed monomial]
\lean{FABL.granularityIndexedMonomial}
\label{FABL.granularityIndexedMonomial}
\leanok
\uses{FABL.Sign, FABL.signValue}
Let $\iota$ be a type, let $S$ be a finite set of indices in $\iota$, and let $x$ assign
a unit sign to each index. The granularity indexed monomial of $x$ over $S$ is the
product of the real sign values of $x_i$ for $i\in S$.
\end{definition}

\begin{definition}[Sign monomial character]
\lean{FABL.granularitySignMonomialChar}
\label{FABL.granularitySignMonomialChar}
\leanok
\uses{FABL.Sign, FABL.granularityIndexedMonomial, FABL.signValue}
Let $\iota$ be a type and let $S$ be a finite subset of $\iota$. The character on the
additive group of functions from $\iota$ to the signs is the map that sends each
function to the product of its coordinate signs over $S$.
\end{definition}

\begin{theorem}[Injectivity of sign monomial characters]
\lean{FABL.granularitySignMonomialChar_injective}
\label{FABL.granularitySignMonomialChar_injective}
\leanok
\uses{FABL.Sign, FABL.granularityIndexedMonomial, FABL.granularitySignMonomialChar, FABL.signValue}
\difficulty{5}
For any type $\iota$, associate to each finite subset $S\subseteq\iota$ the character on
sign-valued functions $x:\iota\to\{\pm1\}$ given by $x\mapsto\prod_{i\in S}x(i)$. This
association is injective: distinct finite subsets define distinct characters.
\end{theorem}

\begin{theorem}[Orthogonality of sign monomials]
\lean{FABL.expect_granularityIndexedMonomial_mul}
\label{FABL.expect_granularityIndexedMonomial_mul}
\leanok
\uses{FABL.Sign, FABL.granularityIndexedMonomial, FABL.granularitySignMonomialChar, FABL.granularitySignMonomialChar_injective}
\difficulty{4}
Let $I$ be a finite set and let $S,T\\subseteq I$. For a uniformly random sign
assignment $x\\in\\{-1,1\\}^I$, let $m_S(x)=\\prod_{i\\in S}x_i$. Then $\\mathbb
E[m_S(x)m_T(x)]$ equals $1$ if $S=T$ and $0$ otherwise.
\end{theorem}

\begin{definition}[Extension by the positive sign]
\lean{FABL.granularityExtendAtOne}
\label{FABL.granularityExtendAtOne}
\leanok
\uses{FABL.Sign}
Let $n$ be a natural number, let $J$ be a finite set of coordinates in
$\{0,\ldots,n-1\}$, and let $y$ assign a sign to each coordinate in $J$. Extend $y$ to
all $n$ coordinates by keeping its values on $J$ and assigning $+1$ outside $J$.
\end{definition}

\begin{definition}[Coordinates of one set lying in another]
\lean{FABL.granularityFrequencyWithin}
\label{FABL.granularityFrequencyWithin}
\leanok
Let $n$ be a natural number and let $J,U\subseteq\{0,\ldots,n-1\}$ be finite. The
selected subset of $J$ consists of exactly those coordinates of $J$ that also belong to
$U$.
\end{definition}

\begin{theorem}[Membership in a restricted frequency set]
\lean{FABL.mem_granularityFrequencyWithin}
\label{FABL.mem_granularityFrequencyWithin}
\leanok
\uses{FABL.granularityFrequencyWithin}
\difficulty{1}
Let $n$ be a natural number, let $J$ and $U$ be sets of coordinates in
$\\mathrm{Fin}(n)$, and let $i\\in J$. Then $i$ belongs to the subset of $J$ consisting
of coordinates that lie in $U$ if and only if the underlying coordinate of $i$ belongs
to $U$.
\end{theorem}

\begin{theorem}[Full frequency set criterion]
\lean{FABL.granularityFrequencyWithin_eq_univ_iff}
\label{FABL.granularityFrequencyWithin_eq_univ_iff}
\leanok
\uses{FABL.granularityFrequencyWithin, FABL.mem_granularityFrequencyWithin}
\difficulty{4}
Let $J$ and $U$ be finite subsets of $\{0,\ldots,n-1\}$. The elements of $J$ that lie in
$U$ comprise all of $J$ if and only if $J\subseteq U$.
\end{theorem}

\begin{theorem}[Monomial on an extension by positive signs]
\lean{FABL.monomial_granularityExtendAtOne}
\label{FABL.monomial_granularityExtendAtOne}
\leanok
\uses{FABL.Sign, FABL.granularityExtendAtOne, FABL.granularityFrequencyWithin, FABL.granularityIndexedMonomial, FABL.mem_granularityFrequencyWithin, FABL.monomial, FABL.signValue}
\difficulty{5}
Let $J$ and $U$ be finite sets of coordinates among $n$, and let $y$ assign a sign to
each element of $J$. Extend $y$ to all $n$ coordinates by assigning $+1$ outside $J$.
The monomial on $U$ at this extension equals the monomial on the coordinates in $U\cap
J$, evaluated using $y$.
\end{theorem}

\begin{definition}[Top coefficient of a sign restriction]
\lean{FABL.granularityRestrictionTopCoeff}
\label{FABL.granularityRestrictionTopCoeff}
\leanok
\uses{FABL.Sign, FABL.granularityExtendAtOne, FABL.granularityIndexedMonomial, FABL.signValue}
Let $n$ be a natural number, let $f$ map sign vectors of length $n$ to signs, and let
$J$ be a set of coordinates among the first $n$ coordinates. The top coefficient of the
restriction of $f$ to $J$, with all coordinates outside $J$ fixed to $+1$, is the
average over all sign assignments $y$ on $J$ of the value of $f$ at that extended
assignment multiplied by the product of the signs assigned by $y$ to every coordinate in
$J$.
\end{definition}

\begin{theorem}[Indicator polynomial fourier sum]
\lean{FABL.indicatorPolynomial_fourier_sum}
\statementsource{Analysis-of-Boolean-Functions-by-Ryan-ODonnell}{
  pdf-ffbe9c7b3306-p021-b002,
  pdf-ffbe9c7b3306-p029-b006
}
\label{FABL.indicatorPolynomial_fourier_sum}
\leanok
\uses{FABL.indicatorPolynomial, FABL.monomial, FABL.signValue}
\difficulty{4}
% dataset-pipeline: stub -- prose pending build_informal_statements.py
Statement pending regeneration.
\end{theorem}

\begin{theorem}[Fourier expansion from interpolation]
\lean{FABL.fourier_expansion_from_interpolation}
\statementsource{lmn-constant-depth-fourier}{pdf-63759bd909d5-p004-b003}
\proofsource{Analysis-of-Boolean-Functions-by-Ryan-ODonnell}{
  pdf-ffbe9c7b3306-p021-b002,
  pdf-ffbe9c7b3306-p022-b002
}
\statementsource{mossel-neeman-gaussian-noise-stability}{pdf-b53c1a5aa27d-p034-b004}
\label{FABL.fourier_expansion_from_interpolation}
\leanok
\uses{FABL.fourierCoeff, FABL.indicatorPolynomial, FABL.indicatorPolynomial_fourier_sum, FABL.monomial, FABL.sum_indicatorPolynomial}
\proofstep
  {FABL.Sign}
  {context}
  {Analysis-of-Boolean-Functions-by-Ryan-ODonnell}
  {pdf-ffbe9c7b3306-p014-b002,
    pdf-ffbe9c7b3306-p014-b003}
\proofstep
  {FABL.signValue}
  {context}
  {Analysis-of-Boolean-Functions-by-Ryan-ODonnell}
  {pdf-ffbe9c7b3306-p014-b002}
\proofstep
  {FABL.signValue_one}
  {context}
  {Analysis-of-Boolean-Functions-by-Ryan-ODonnell}
  {pdf-ffbe9c7b3306-p014-b002}
\proofstep
  {FABL.signValue_neg_one}
  {context}
  {Analysis-of-Boolean-Functions-by-Ryan-ODonnell}
  {pdf-ffbe9c7b3306-p014-b002}
\proofstep
  {FABL.indicatorPolynomial}
  {direct}
  {Analysis-of-Boolean-Functions-by-Ryan-ODonnell}
  {pdf-ffbe9c7b3306-p021-b002}
\proofstep
  {FABL.indicatorPolynomial_eq_ite}
  {direct}
  {Analysis-of-Boolean-Functions-by-Ryan-ODonnell}
  {pdf-ffbe9c7b3306-p021-b002}
\proofstep
  {FABL.sum_indicatorPolynomial}
  {direct}
  {Analysis-of-Boolean-Functions-by-Ryan-ODonnell}
  {pdf-ffbe9c7b3306-p021-b002}
\proofstep
  {FABL.monomial}
  {direct}
  {Analysis-of-Boolean-Functions-by-Ryan-ODonnell}
  {pdf-ffbe9c7b3306-p022-b001,
    pdf-ffbe9c7b3306-p022-b004}
\proofstep
  {FABL.fourierCoeff}
  {direct}
  {Analysis-of-Boolean-Functions-by-Ryan-ODonnell}
  {pdf-ffbe9c7b3306-p025-b004}
\proofstep
  {FABL.fourierSupport}
  {context}
  {Analysis-of-Boolean-Functions-by-Ryan-ODonnell}
  {pdf-ffbe9c7b3306-p022-b002,
    pdf-ffbe9c7b3306-p025-b004}
\proofstep
  {FABL.mem_fourierSupport}
  {context}
  {Analysis-of-Boolean-Functions-by-Ryan-ODonnell}
  {pdf-ffbe9c7b3306-p022-b002}
\proofstep
  {FABL.granularityFrequencyWithin}
  {context}
  {Analysis-of-Boolean-Functions-by-Ryan-ODonnell}
  {pdf-ffbe9c7b3306-p075-b002,
    pdf-ffbe9c7b3306-p075-b005}
\proofstep
  {FABL.mem_granularityFrequencyWithin}
  {context}
  {Analysis-of-Boolean-Functions-by-Ryan-ODonnell}
  {pdf-ffbe9c7b3306-p075-b002,
    pdf-ffbe9c7b3306-p075-b005}
\proofstep
  {FABL.indicatorPolynomial_fourier_sum}
  {context}
  {Analysis-of-Boolean-Functions-by-Ryan-ODonnell}
  {pdf-ffbe9c7b3306-p021-b002,
    pdf-ffbe9c7b3306-p022-b002,
    pdf-ffbe9c7b3306-p023-b004}
\difficulty{4}

% dataset-pipeline: stub -- prose pending build_informal_statements.py
Statement pending regeneration.
\end{theorem}

\begin{theorem}[Granularity restriction top coeff eq sum supersets]
\lean{FABL.granularityRestrictionTopCoeff_eq_sum_supersets}
\statementsource{Analysis-of-Boolean-Functions-by-Ryan-ODonnell}{
  pdf-ffbe9c7b3306-p076-b003,
  pdf-ffbe9c7b3306-p076-b004
}
\label{FABL.granularityRestrictionTopCoeff_eq_sum_supersets}
\leanok
\uses{FABL.Sign, FABL.expect_granularityIndexedMonomial_mul, FABL.fourierCoeff, FABL.fourier_expansion_from_interpolation, FABL.granularityExtendAtOne, FABL.granularityFrequencyWithin, FABL.granularityFrequencyWithin_eq_univ_iff, FABL.granularityIndexedMonomial, FABL.granularityRestrictionTopCoeff, FABL.monomial, FABL.monomial_granularityExtendAtOne, FABL.signValue}
\difficulty{4}
% dataset-pipeline: stub -- prose pending build_informal_statements.py
Statement pending regeneration.
\end{theorem}

\begin{theorem}[Mean of a sign-valued function]
\lean{FABL.expect_sign_values_eq_int_mul_two_inv_pow}
\label{FABL.expect_sign_values_eq_int_mul_two_inv_pow}
\leanok
\difficulty{4}
Let $\Omega$ be a finite nonempty set, let $g:\Omega\to\mathbb{R}$ take only the values
$-1$ and $1$, and let $m\ge 1$ be an integer such that $|\Omega|=2^m$. There is an
integer $z$ for which the uniform mean of $g$ is $z\,2(2^{-1})^m$.
\end{theorem}

\begin{theorem}[Integer scaling for a nonempty restriction set]
\lean{FABL.granularityRestrictionTopCoeff_eq_int_mul_two_inv_pow_of_nonempty}
\label{FABL.granularityRestrictionTopCoeff_eq_int_mul_two_inv_pow_of_nonempty}
\leanok
\uses{FABL.Sign, FABL.expect_sign_values_eq_int_mul_two_inv_pow, FABL.granularityExtendAtOne, FABL.granularityIndexedMonomial, FABL.granularityRestrictionTopCoeff, FABL.signValue, FABL.signValue_eq_neg_one_or_one}
\difficulty{3}
Let $f:\{\pm1\}^n\to\{\pm1\}$ and let $J\subseteq\{0,\ldots,n-1\}$ be a nonempty finite
set. There is an integer $z$ such that the top restriction coefficient of $f$ on $J$
equals $z\,2(1/2)^{|J|}$.
\end{theorem}

\begin{theorem}[Integer scaling of the top restriction coefficient]
\lean{FABL.granularityRestrictionTopCoeff_eq_int_mul_global_scale}
\proofsource{Analysis-of-Boolean-Functions-by-Ryan-ODonnell}{pdf-ffbe9c7b3306-p036-b004}
\label{FABL.granularityRestrictionTopCoeff_eq_int_mul_global_scale}
\leanok
\uses{FABL.Sign, FABL.granularityExtendAtOne, FABL.granularityIndexedMonomial, FABL.granularityRestrictionTopCoeff, FABL.granularityRestrictionTopCoeff_eq_int_mul_two_inv_pow_of_nonempty, FABL.signValue, FABL.signValue_eq_neg_one_or_one}
\proofstep
  {FABL.Sign}
  {context}
  {Analysis-of-Boolean-Functions-by-Ryan-ODonnell}
  {pdf-ffbe9c7b3306-p020-b002,
    pdf-ffbe9c7b3306-p022-b004}
\proofstep
  {FABL.signValue}
  {context}
  {Analysis-of-Boolean-Functions-by-Ryan-ODonnell}
  {pdf-ffbe9c7b3306-p022-b004}
\proofstep
  {FABL.signValue_one}
  {context}
  {Analysis-of-Boolean-Functions-by-Ryan-ODonnell}
  {pdf-ffbe9c7b3306-p022-b004}
\proofstep
  {FABL.signValue_neg_one}
  {context}
  {Analysis-of-Boolean-Functions-by-Ryan-ODonnell}
  {pdf-ffbe9c7b3306-p022-b004}
\proofstep
  {FABL.signValue_eq_neg_one_or_one}
  {context}
  {Analysis-of-Boolean-Functions-by-Ryan-ODonnell}
  {pdf-ffbe9c7b3306-p020-b002,
    pdf-ffbe9c7b3306-p022-b004}
\proofstep
  {FABL.monomial}
  {direct}
  {Analysis-of-Boolean-Functions-by-Ryan-ODonnell}
  {pdf-ffbe9c7b3306-p022-b001,
    pdf-ffbe9c7b3306-p022-b004}
\proofstep
  {FABL.fourierCoeff}
  {direct}
  {Analysis-of-Boolean-Functions-by-Ryan-ODonnell}
  {pdf-ffbe9c7b3306-p025-b004}
\proofstep
  {FABL.fourierSupport}
  {context}
  {Analysis-of-Boolean-Functions-by-Ryan-ODonnell}
  {pdf-ffbe9c7b3306-p022-b002,
    pdf-ffbe9c7b3306-p027-b008}
\proofstep
  {FABL.mem_fourierSupport}
  {context}
  {Analysis-of-Boolean-Functions-by-Ryan-ODonnell}
  {pdf-ffbe9c7b3306-p022-b002,
    pdf-ffbe9c7b3306-p027-b008}
\proofstep
  {FABL.granularityIndexedMonomial}
  {context}
  {Analysis-of-Boolean-Functions-by-Ryan-ODonnell}
  {pdf-ffbe9c7b3306-p022-b004}
\proofstep
  {FABL.granularityExtendAtOne}
  {context}
  {Analysis-of-Boolean-Functions-by-Ryan-ODonnell}
  {pdf-ffbe9c7b3306-p075-b002}
\proofstep
  {FABL.granularityFrequencyWithin}
  {context}
  {Analysis-of-Boolean-Functions-by-Ryan-ODonnell}
  {pdf-ffbe9c7b3306-p076-b001}
\proofstep
  {FABL.mem_granularityFrequencyWithin}
  {context}
  {Analysis-of-Boolean-Functions-by-Ryan-ODonnell}
  {pdf-ffbe9c7b3306-p076-b001}
\proofstep
  {FABL.granularityRestrictionTopCoeff}
  {context}
  {Analysis-of-Boolean-Functions-by-Ryan-ODonnell}
  {pdf-ffbe9c7b3306-p075-b002,
    pdf-ffbe9c7b3306-p076-b001,
    pdf-ffbe9c7b3306-p076-b003}
\proofstep
  {FABL.expect_sign_values_eq_int_mul_two_inv_pow}
  {direct}
  {Analysis-of-Boolean-Functions-by-Ryan-ODonnell}
  {pdf-ffbe9c7b3306-p036-b004}
\proofstep
  {FABL.granularityRestrictionTopCoeff_eq_int_mul_two_inv_pow_of_nonempty}
  {context}
  {Analysis-of-Boolean-Functions-by-Ryan-ODonnell}
  {pdf-ffbe9c7b3306-p075-b002,
    pdf-ffbe9c7b3306-p076-b001,
    pdf-ffbe9c7b3306-p036-b004}
\difficulty{5}
Let $f:\{\pm1\}^n\to\{\pm1\}$, let $J\subseteq\{0,\ldots,n-1\}$ be finite, and let
$k\ge1$ satisfy $|J|\le k$. There is an integer $z$ such that the top restriction
coefficient of $f$ on $J$ equals $z\,2(1/2)^k$.
\end{theorem}

\begin{theorem}[Fourier coeff eq int mul global scale of card le]
\lean{FABL.fourierCoeff_eq_int_mul_global_scale_of_card_le}
\statementsource{Analysis-of-Boolean-Functions-by-Ryan-ODonnell}{pdf-ffbe9c7b3306-p036-b004}
\label{FABL.fourierCoeff_eq_int_mul_global_scale_of_card_le}
\leanok
\uses{FABL.Sign, FABL.fourierCoeff, FABL.fourierDegree, FABL.fourierDegree_le_iff, FABL.granularityRestrictionTopCoeff, FABL.granularityRestrictionTopCoeff_eq_int_mul_global_scale, FABL.granularityRestrictionTopCoeff_eq_sum_supersets, FABL.signValue}
\difficulty{6}
% dataset-pipeline: stub -- prose pending build_informal_statements.py
Statement pending regeneration.
\end{theorem}

\begin{theorem}[Granularity from Bounded Fourier Degree]
\lean{FABL.isFourierGranular_signValue_of_fourierDegree_le}
\proofsource{Analysis-of-Boolean-Functions-by-Ryan-ODonnell}{
  pdf-ffbe9c7b3306-p036-b004,
  pdf-ffbe9c7b3306-p072-b003,
  pdf-ffbe9c7b3306-p089-b003
}
\label{FABL.isFourierGranular_signValue_of_fourierDegree_le}
\leanok
\uses{FABL.IsFourierGranular, FABL.Sign, FABL.fourierCoeff_eq_int_mul_global_scale_of_card_le, FABL.fourierDegree, FABL.fourierDegree_le_iff, FABL.signValue}
\proofstep
  {FABL.Sign}
  {context}
  {Analysis-of-Boolean-Functions-by-Ryan-ODonnell}
  {pdf-ffbe9c7b3306-p020-b001,
    pdf-ffbe9c7b3306-p014-b002}
\proofstep
  {FABL.signValue}
  {context}
  {Analysis-of-Boolean-Functions-by-Ryan-ODonnell}
  {pdf-ffbe9c7b3306-p014-b002,
    pdf-ffbe9c7b3306-p022-b004}
\proofstep
  {FABL.signValue_one}
  {context}
  {Analysis-of-Boolean-Functions-by-Ryan-ODonnell}
  {pdf-ffbe9c7b3306-p014-b002}
\proofstep
  {FABL.signValue_neg_one}
  {context}
  {Analysis-of-Boolean-Functions-by-Ryan-ODonnell}
  {pdf-ffbe9c7b3306-p014-b002}
\proofstep
  {FABL.signValue_eq_neg_one_or_one}
  {context}
  {Analysis-of-Boolean-Functions-by-Ryan-ODonnell}
  {pdf-ffbe9c7b3306-p020-b001}
\proofstep
  {FABL.indicatorPolynomial}
  {context}
  {Analysis-of-Boolean-Functions-by-Ryan-ODonnell}
  {pdf-ffbe9c7b3306-p021-b002}
\proofstep
  {FABL.indicatorPolynomial_eq_ite}
  {direct}
  {Analysis-of-Boolean-Functions-by-Ryan-ODonnell}
  {pdf-ffbe9c7b3306-p021-b002}
\proofstep
  {FABL.sum_indicatorPolynomial}
  {direct}
  {Analysis-of-Boolean-Functions-by-Ryan-ODonnell}
  {pdf-ffbe9c7b3306-p021-b002,
    pdf-ffbe9c7b3306-p021-b004}
\proofstep
  {FABL.monomial}
  {direct}
  {Analysis-of-Boolean-Functions-by-Ryan-ODonnell}
  {pdf-ffbe9c7b3306-p022-b001,
    pdf-ffbe9c7b3306-p022-b004}
\proofstep
  {FABL.fourierCoeff}
  {direct}
  {Analysis-of-Boolean-Functions-by-Ryan-ODonnell}
  {pdf-ffbe9c7b3306-p025-b004}
\proofstep
  {FABL.fourierSupport}
  {context}
  {Analysis-of-Boolean-Functions-by-Ryan-ODonnell}
  {pdf-ffbe9c7b3306-p022-b002}
\proofstep
  {FABL.fourierDegree}
  {context}
  {Analysis-of-Boolean-Functions-by-Ryan-ODonnell}
  {pdf-ffbe9c7b3306-p036-b002}
\proofstep
  {FABL.mem_fourierSupport}
  {context}
  {Analysis-of-Boolean-Functions-by-Ryan-ODonnell}
  {pdf-ffbe9c7b3306-p022-b002}
\proofstep
  {FABL.fourierDegree_le_iff}
  {context}
  {Analysis-of-Boolean-Functions-by-Ryan-ODonnell}
  {pdf-ffbe9c7b3306-p036-b002,
    pdf-ffbe9c7b3306-p036-b003}
\proofstep
  {FABL.IsFourierGranular}
  {direct}
  {Analysis-of-Boolean-Functions-by-Ryan-ODonnell}
  {pdf-ffbe9c7b3306-p072-b003}
\proofstep
  {FABL.granularityIndexedMonomial}
  {context}
  {Analysis-of-Boolean-Functions-by-Ryan-ODonnell}
  {pdf-ffbe9c7b3306-p014-b002,
    pdf-ffbe9c7b3306-p024-b005}
\proofstep
  {FABL.granularitySignMonomialChar}
  {context}
  {Analysis-of-Boolean-Functions-by-Ryan-ODonnell}
  {pdf-ffbe9c7b3306-p023-b001,
    pdf-ffbe9c7b3306-p024-b003}
\proofstep
  {FABL.granularitySignMonomialChar_injective}
  {context}
  {Analysis-of-Boolean-Functions-by-Ryan-ODonnell}
  {pdf-ffbe9c7b3306-p024-b003}
\proofstep
  {FABL.expect_granularityIndexedMonomial_mul}
  {context}
  {Analysis-of-Boolean-Functions-by-Ryan-ODonnell}
  {pdf-ffbe9c7b3306-p024-b003,
    pdf-ffbe9c7b3306-p024-b007}
\proofstep
  {FABL.granularityExtendAtOne}
  {context}
  {Analysis-of-Boolean-Functions-by-Ryan-ODonnell}
  {pdf-ffbe9c7b3306-p075-b002}
\proofstep
  {FABL.granularityFrequencyWithin}
  {context}
  {Analysis-of-Boolean-Functions-by-Ryan-ODonnell}
  {pdf-ffbe9c7b3306-p076-b001}
\proofstep
  {FABL.mem_granularityFrequencyWithin}
  {context}
  {Analysis-of-Boolean-Functions-by-Ryan-ODonnell}
  {pdf-ffbe9c7b3306-p076-b001}
\proofstep
  {FABL.granularityFrequencyWithin_eq_univ_iff}
  {context}
  {Analysis-of-Boolean-Functions-by-Ryan-ODonnell}
  {pdf-ffbe9c7b3306-p076-b001}
\proofstep
  {FABL.monomial_granularityExtendAtOne}
  {context}
  {Analysis-of-Boolean-Functions-by-Ryan-ODonnell}
  {pdf-ffbe9c7b3306-p075-b002,
    pdf-ffbe9c7b3306-p076-b001}
\proofstep
  {FABL.granularityRestrictionTopCoeff}
  {context}
  {Analysis-of-Boolean-Functions-by-Ryan-ODonnell}
  {pdf-ffbe9c7b3306-p075-b002,
    pdf-ffbe9c7b3306-p025-b004}
\proofstep
  {FABL.indicatorPolynomial_fourier_sum}
  {direct}
  {Analysis-of-Boolean-Functions-by-Ryan-ODonnell}
  {pdf-ffbe9c7b3306-p021-b002,
    pdf-ffbe9c7b3306-p022-b002}
\proofstep
  {FABL.fourier_expansion_from_interpolation}
  {direct}
  {Analysis-of-Boolean-Functions-by-Ryan-ODonnell}
  {pdf-ffbe9c7b3306-p021-b002,
    pdf-ffbe9c7b3306-p022-b002}
\proofstep
  {FABL.granularityRestrictionTopCoeff_eq_sum_supersets}
  {context}
  {Analysis-of-Boolean-Functions-by-Ryan-ODonnell}
  {pdf-ffbe9c7b3306-p075-b002,
    pdf-ffbe9c7b3306-p022-b002,
    pdf-ffbe9c7b3306-p024-b003}
\proofstep
  {FABL.expect_sign_values_eq_int_mul_two_inv_pow}
  {direct}
  {Analysis-of-Boolean-Functions-by-Ryan-ODonnell}
  {pdf-ffbe9c7b3306-p036-b004}
\proofstep
  {FABL.granularityRestrictionTopCoeff_eq_int_mul_two_inv_pow_of_nonempty}
  {context}
  {Analysis-of-Boolean-Functions-by-Ryan-ODonnell}
  {pdf-ffbe9c7b3306-p036-b004}
\proofstep
  {FABL.granularityRestrictionTopCoeff_eq_int_mul_global_scale}
  {context}
  {Analysis-of-Boolean-Functions-by-Ryan-ODonnell}
  {pdf-ffbe9c7b3306-p036-b004}
\proofstep
  {FABL.fourierCoeff_eq_int_mul_global_scale_of_card_le}
  {direct}
  {Analysis-of-Boolean-Functions-by-Ryan-ODonnell}
  {pdf-ffbe9c7b3306-p036-b004}
\difficulty{4}
Let $f:\{\pm1\}^n\to\{\pm1\}$, and let $k\ge 1$ be an integer such that every nonzero
Fourier coefficient of $f$ is indexed by a set of size at most $k$. Then each Fourier
coefficient of $f$ is an integer multiple of $2^{1-k}$.
\end{theorem}

\begin{definition}[Multilinear polynomial on the sign cube]
\lean{FABL.multilinearPolynomial}
\statementsource{lmn-constant-depth-fourier}{pdf-63759bd909d5-p004-b006}
\statementsource{Analysis-of-Boolean-Functions-by-Ryan-ODonnell}{
  pdf-ffbe9c7b3306-p020-b005,
  pdf-ffbe9c7b3306-p021-b004,
  pdf-ffbe9c7b3306-p149-b001,
  pdf-ffbe9c7b3306-p149-b002,
  pdf-ffbe9c7b3306-p215-b009,
  pdf-ffbe9c7b3306-p215-b010,
  pdf-ffbe9c7b3306-p216-b001
}
\statementsource{mossel-neeman-gaussian-noise-stability}{pdf-b53c1a5aa27d-p034-b006}
\statementsource{goldreich-approximation-method}{
  pdf-988dc9f0ae91-p009-b001,
  pdf-988dc9f0ae91-p010-b001,
  pdf-988dc9f0ae91-p012-b003,
  pdf-988dc9f0ae91-p013-b001
}
\label{FABL.multilinearPolynomial}
\leanok
\uses{FABL.monomial}
Given real coefficients $a(S)$ indexed by subsets $S$ of the coordinates, the associated
multilinear polynomial on the sign cube is $P(x)=\sum_S a(S)\prod_{i\in S}x_i$.
\end{definition}

\begin{theorem}[Fourier expansion unique]
\lean{FABL.fourier_expansion_unique}
\proofsource{lmn-constant-depth-fourier}{pdf-63759bd909d5-p004-b003}
\statementsource{Analysis-of-Boolean-Functions-by-Ryan-ODonnell}{
  pdf-ffbe9c7b3306-p022-b002,
  pdf-ffbe9c7b3306-p023-b006
}
\label{FABL.fourier_expansion_unique}
\leanok
\uses{FABL.expect_monomial_mul, FABL.fourierCoeff, FABL.fourier_expansion_from_interpolation, FABL.monomial, FABL.multilinearPolynomial}
\proofstep
  {FABL.Sign}
  {context}
  {lmn-constant-depth-fourier}
  {pdf-63759bd909d5-p004-b002}
\proofstep
  {FABL.signValue}
  {context}
  {lmn-constant-depth-fourier}
  {pdf-63759bd909d5-p004-b002}
\proofstep
  {FABL.signValue_one}
  {context}
  {lmn-constant-depth-fourier}
  {pdf-63759bd909d5-p004-b002}
\proofstep
  {FABL.signValue_neg_one}
  {context}
  {lmn-constant-depth-fourier}
  {pdf-63759bd909d5-p004-b002}
\proofstep
  {FABL.indicatorPolynomial}
  {context}
  {lmn-constant-depth-fourier}
  {pdf-63759bd909d5-p004-b003}
\proofstep
  {FABL.indicatorPolynomial_eq_ite}
  {context}
  {lmn-constant-depth-fourier}
  {pdf-63759bd909d5-p004-b003}
\proofstep
  {FABL.sum_indicatorPolynomial}
  {context}
  {lmn-constant-depth-fourier}
  {pdf-63759bd909d5-p004-b003}
\proofstep
  {FABL.monomial}
  {context}
  {lmn-constant-depth-fourier}
  {pdf-63759bd909d5-p004-b002}
\proofstep
  {FABL.signMonomialChar}
  {context}
  {lmn-constant-depth-fourier}
  {pdf-63759bd909d5-p004-b002,
    pdf-63759bd909d5-p004-b003}
\proofstep
  {FABL.signMonomialChar_injective}
  {context}
  {lmn-constant-depth-fourier}
  {pdf-63759bd909d5-p004-b003}
\proofstep
  {FABL.expect_monomial_mul}
  {direct}
  {lmn-constant-depth-fourier}
  {pdf-63759bd909d5-p004-b003}
\proofstep
  {FABL.fourierCoeff}
  {direct}
  {lmn-constant-depth-fourier}
  {pdf-63759bd909d5-p004-b003}
\proofstep
  {FABL.fourierSupport}
  {context}
  {lmn-constant-depth-fourier}
  {pdf-63759bd909d5-p004-b003}
\proofstep
  {FABL.mem_fourierSupport}
  {context}
  {lmn-constant-depth-fourier}
  {pdf-63759bd909d5-p004-b003}
\proofstep
  {FABL.granularityFrequencyWithin}
  {context}
  {lmn-constant-depth-fourier}
  {pdf-63759bd909d5-p004-b003}
\proofstep
  {FABL.mem_granularityFrequencyWithin}
  {context}
  {lmn-constant-depth-fourier}
  {pdf-63759bd909d5-p004-b003}
\proofstep
  {FABL.indicatorPolynomial_fourier_sum}
  {context}
  {lmn-constant-depth-fourier}
  {pdf-63759bd909d5-p004-b003}
\proofstep
  {FABL.fourier_expansion_from_interpolation}
  {direct}
  {lmn-constant-depth-fourier}
  {pdf-63759bd909d5-p004-b003}
\proofstep
  {FABL.multilinearPolynomial}
  {context}
  {lmn-constant-depth-fourier}
  {pdf-63759bd909d5-p004-b006,
    pdf-63759bd909d5-p011-b005}
\difficulty{4}

% dataset-pipeline: stub -- prose pending build_informal_statements.py
Statement pending regeneration.
\end{theorem}

\begin{theorem}[Fourier expansion]
\lean{FABL.fourier_expansion}
\statementsource{lmn-constant-depth-fourier}{pdf-63759bd909d5-p004-b003}
\proofsource{Analysis-of-Boolean-Functions-by-Ryan-ODonnell}{
  pdf-ffbe9c7b3306-p022-b002,
  pdf-ffbe9c7b3306-p209-b005
}
\statementsource{mossel-neeman-gaussian-noise-stability}{pdf-b53c1a5aa27d-p034-b004}
\label{FABL.fourier_expansion}
\leanok
\uses{FABL.fourierCoeff, FABL.fourier_expansion_unique, FABL.monomial, FABL.multilinearPolynomial}
\difficulty{1}
% dataset-pipeline: stub -- prose pending build_informal_statements.py
Statement pending regeneration.
\end{theorem}
